\clearpage
\SetRegion{第三章 · 主体影响}{violet}{}
\pdfbookmark[1]{本章导读：由榜单进入持续供给}{ch3-guide}
\kicker{GOSR / 全球开源影响力排行榜}
\PageTitle{由榜单进入持续供给}{六类主体，一个共同问题：影响由哪些资源与协作形成？}
第二章观察参与和资源的分布，本章将视野推进到具体主体。先由完整榜单定位，再比较结构与变化，最后进入项目的交付、复用和维护。

\noindent\includegraphics[width=\linewidth]{figures/new/chapter3_navigation.pdf}\par
\vspace{3mm}
\RunTitle{从完整榜单，读到完整项目}
3.4以OpenHarmony和PyTorch展开分析，并保留两张A3项目画像。前者说明多仓系统的资源组织与跨平台观察，后者连接贡献接续、版本交付与维护责任；二者不合成统一健康分。

\vfill
\ReportNote{主统计期为2025-07-01至2026-06-30。原榜缺少时间字段的图明确标注“统计窗口待核”；历史标签、2026年6月模型月榜及期后快照分别注明。OpenRank与ModelRank各自解释，不直接比较，也不替代质量、治理或公共成效判断。}
